% ============================================================================
%  SECCIÓN 3.22: CONFIGURACIÓN DEL TEMA VISUAL
% ============================================================================

\section{Configuraci\'on del tema visual}

La configuraci\'on del tema visual es un requisito fundamental del avance,
pues asegura que todas las pantallas mantengan una identidad visual coherente
\citep{nielsen1994}. En esta secci\'on se documenta la paleta de colores, la
tipograf\'ia, la escala de espaciado y los estilos de componentes que
conforman el sistema de dise\~no de la aplicaci\'on \nombreApp{}.

\subsection{Paleta de colores}

La paleta de colores se define en el archivo \texttt{src/theme/colors.ts}
y sigue una estructura de tokens de dise\~no inspirada en la escala de
colores de Tailwind CSS. Cada color se asocia a un nombre sem\'antico que
describe su funci\'on en la interfaz.

\begin{table}[H]
  \centering
  \caption{Paleta de colores de la aplicaci\'on \nombreApp{}}
  \contenidoConFuente{%
    \begin{tablaAPA}
      \begin{tabular}{@{}p{2.2cm}p{2.6cm}p{2.2cm}p{8.6cm}@{}}
        \toprule
        \textbf{Categor\'ia} & \textbf{Token} & \textbf{Valor} &
        \textbf{Uso} \\\\
        \midrule
        \textbf{Primario} & \textit{primary} & \textit{\#2563EB} & Color
        principal de la marca: botones, enlaces, \'iconos activos. \\\\
        \addlinespace
        Primario & \textit{primaryDark} & \textit{\#1D4ED8} & Fondo de la
        parte superior del splash. \\\\
        \addlinespace
        Primario & \textit{primaryLight} & \textit{\#93C5FD} & Texto tagline
        del splash, elementos secundarios. \\\\
        \addlinespace
        \textbf{Secundario} & \textit{secondary} & \textit{\#10B981} & Color
        de \'exito, confirmaciones positivas. \\\\
        \addlinespace
        \textbf{Superficie} & \textit{background} & \textit{\#F8FAFC} &
        Fondo general de todas las pantallas. \\\\
        \addlinespace
        Superficie & \textit{surface} & \textit{\#FFFFFF} & Fondo de
        tarjetas, campos de formulario y botones. \\\\
        \addlinespace
        \textbf{Texto} & \textit{text} & \textit{\#0F172A} & Color principal
        del texto (t\'itulos, contenido). \\\\
        \addlinespace
        Texto & \textit{textSecondary} & \textit{\#64748B} & Texto
        secundario, subt\'itulos y placeholders. \\\\
        \addlinespace
        \textbf{Estado} & \textit{error} & \textit{\#EF4444} & Errores de
        validaci\'on y mensajes de error del servidor. \\\\
        \addlinespace
        Estado & \textit{success} & \textit{\#22C55E} & Mensajes de
        \'exito (registro exitoso). \\\\
        \addlinespace
        Estado & \textit{warning} & \textit{\#F59E0B} & Advertencias
        (reservado para avances futuros). \\\\
        \bottomrule
      \end{tabular}
    \end{tablaAPA}%
  }{Elaboraci\'on propia en base al archivo \texttt{colors.ts} del proyecto.}
\end{table}

\subsection{Escalas de tipograf\'ia y espaciado}

La aplicaci\'on utiliza la tipograf\'ia del sistema (\texttt{fontFamily}
por defecto de React Native), que var\'ia seg\'un la plataforma: San Francisco
en iOS y Roboto en Android. Se define una escala de tama\~nos de fuente que
va desde 12\,pt (\texttt{xs}) hasta 32\,pt (\texttt{title}).

\begin{table}[H]
  \centering
  \caption{Escalas de tipograf\'ia y espaciado del tema}
  \contenidoConFuente{%
    \begin{tablaAPA}
      \begin{tabular}{@{}p{3.2cm}p{3.6cm}p{8.8cm}@{}}
        \toprule
        \textbf{Token} & \textbf{Valor} & \textbf{Uso} \\\\
        \midrule
        \textit{fontSize.xs} & 12\,pt & Mensajes de error y pie de figura. \\\\
        \addlinespace
        \textit{fontSize.sm} & 14\,pt & Etiquetas de campo y texto secundario. \\\\
        \addlinespace
        \textit{fontSize.md} & 16\,pt & Texto de cuerpo y contenido principal. \\\\
        \addlinespace
        \textit{fontSize.lg} & 18\,pt & Texto de botones. \\\\
        \addlinespace
        \textit{fontSize.title} & 32\,pt & T\'itulos de pantalla. \\\\
        \addlinespace
        \textit{spacing.md} & 16\,pt & Espaciado est\'andar entre
        elementos. \\\\
        \addlinespace
        \textit{borderRadius.md} & 8\,pt & Radio de bordes de campos y
        botones. \\\\
        \bottomrule
      \end{tabular}
    \end{tablaAPA}%
  }{Elaboraci\'on propia en base al archivo \texttt{spacing.ts} del proyecto.}
\end{table}

\subsection{Identidad visual coherente}

Todas las pantallas de la aplicaci\'on utilizan la misma paleta de colores y
las mismas escalas de espaciado y tipograf\'ia, lo que garantiza que la
aplicaci\'on se vea como un solo producto y no como pantallas dise\~nadas por
separado. Los componentes de interfaz (botones, campos de texto, tarjetas)
siguen patrones visuales consistentes: botones primarios con fondo de color
\texttt{primary} y texto blanco, campos de formulario con fondo \texttt{surface}
y borde \texttt{border}, y mensajes de error con fondo \texttt{errorLight}.
